\capitulo{7}{Conclusiones y Líneas de trabajo futuras}

El uso de herramientas de simulación en la enseñanza de algoritmos complejos ha demostrado ser un recurso didáctico eficaz para mejorar la comprensión y la retención del conocimiento. Estas herramientas no solo permiten a los docentes ilustrar conceptos abstractos mediante ejemplos concretos, sino que también ofrecen a los estudiantes la oportunidad de reforzar lo aprendido en clase de manera autónoma, a su propio ritmo y tantas veces como sea necesario. Además, facilitan la adquisición independiente de conocimientos al posibilitar la interacción con conjuntos de datos personalizados.

En particular, la herramienta didáctica presentada en este trabajo ha sido diseñada para mejorar la enseñanza y el aprendizaje del proceso de construcción de árboles de decisión mediante el algoritmo C4.5. La posibilidad de visualizar e interactuar con cada una de las etapas del algoritmo contribuye significativamente a la comprensión de conceptos clave, como la entropía, la ganancia de información, la información de división, la razón de ganancia y la selección de umbrales para atributos numéricos.

Asimismo, la herramienta permite comprender cómo C4.5 selecciona los atributos más adecuados para dividir los datos, evitando el sesgo que puede producirse al utilizar exclusivamente la ganancia de información. También facilita el estudio de la generación de árboles con múltiples clases, el tratamiento de atributos numéricos y la aplicación de técnicas de poda para reducir la complejidad del modelo y limitar el sobreajuste.

Al tratarse de una herramienta de código abierto y disponible en línea, se favorece su accesibilidad y aplicabilidad en distintos entornos educativos. Esto permite que tanto docentes como estudiantes puedan utilizarla como recurso de apoyo durante el aprendizaje de minería de datos, aprendizaje automático e inteligencia artificial, promoviendo una metodología más visual, práctica e interactiva.

Aunque se ha realizado una evaluación preliminar de la herramienta, todavía es necesario valorar su efectividad directa en el aprendizaje de los estudiantes, así como analizar con mayor profundidad su experiencia de usuario. Por ello, se plantea como línea de trabajo futuro formalizar su uso dentro de actividades de aprendizaje en asignaturas relacionadas con computación, inteligencia artificial y minería de datos.

Sería especialmente interesante estudiar el impacto de la herramienta en aspectos como la retención de conocimientos, el tiempo necesario para comprender el algoritmo, la capacidad de interpretar árboles de decisión y la mejora de las habilidades prácticas de los estudiantes. Del mismo modo, podrían realizarse estudios longitudinales para medir su impacto educativo a largo plazo y comparar sus resultados con los obtenidos mediante metodologías tradicionales de enseñanza.

Aunque la herramienta permite explicar los elementos principales del algoritmo C4.5, existen diversas posibilidades de ampliación. Una de ellas consistiría en incorporar un mayor número de opciones de configuración, como distintos criterios de poda, valores mínimos de casos por hoja o condiciones de parada ajustables por el usuario. También podría ampliarse la información visual proporcionada durante la ejecución, mostrando de forma más detallada el cálculo de las probabilidades, la distribución de clases en cada nodo y el efecto de la poda sobre la estructura final del árbol.

Otra línea de trabajo futuro sería incorporar algoritmos adicionales de árboles de decisión, como CART (\textit{Classification and Regression Trees}), para permitir la comparación entre distintos criterios de división, como el índice de Gini y la razón de ganancia utilizada por C4.5. Asimismo, podría incluirse la construcción de árboles de regresión, lo que permitiría extender la herramienta a problemas en los que la variable objetivo fuese numérica.

Finalmente, también sería de interés incorporar métodos basados en conjuntos de árboles, como los bosques aleatorios (\textit{Random Forests}) o los algoritmos de \textit{boosting}. Estas ampliaciones permitirían mostrar cómo varios árboles pueden combinarse para mejorar la capacidad predictiva de un modelo y ofrecerían una visión más completa de las técnicas de aprendizaje automático basadas en árboles de decisión.

